%\documentclass[tikz,border=20pt]{standalone}
%\usepackage{amsmath, amssymb}
%\usetikzlibrary{matrix, positioning, decorations.pathreplacing, shadings, shadows, arrows.meta, calc}
%
%\begin{document}
%	\begin{tikzpicture}[
%		>=Stealth,
%		node distance=1.5cm,
%		% Define distinctive colors
%		colorMatrixBg/.style={cyan!10, draw=cyan!60, thick},
%		colorFlattenBg/.style={purple!15, draw=purple!60, thick},
%		colorChoicesBg/.style={orange!20, draw=orange!70, thick},
%		colorFinalBg/.style={yellow!25, draw=yellow!80, thick, drop shadow},
%		% Modern styles for nodes
%		entryNode/.style={draw, minimum size=1.8em, inner sep=0pt, font=\small},
%		choiceButton/.style={circle, minimum size=2.5em, colorChoicesBg, font=\bfseries\large, drop shadow={opacity=0.2}},
%		resultPlaque/.style={rectangle, rounded corners=8pt, colorFinalBg, inner sep=15pt, align=center},
%		labelStyle/.style={font=\sffamily\bfseries\color{darkgray}, align=center},
%		arrowFlow/.style={->, thick, shorten >=2pt, shorten <=2pt},
%		% Corrected brace style (swaps points to avoid mirror)
%		revbrace/.style={decorate, decoration={brace, amplitude=8pt, raise=6pt}}
%		]
%		
%		% ==============================================================================
%		% PHASE 1: THE INPUT MATRIX (Cyan)
%		% ==============================================================================
%		\node[labelStyle] (title1) {$A \in \mathbb{F}_q^{m \times n}$};
%		
%%		\matrix (M) [matrix of math nodes, below=0.5cm of title1,
%%		left delimiter=(, right delimiter=),
%%		nodes={entryNode, fill=cyan!5},
%%		column sep=0.5em, row sep=0.5em,
%%		%fill=cyan!2, 
%%		%draw=cyan!40, 
%%		rounded corners=3pt, 
%%		inner sep=10pt
%%		%, drop shadow={opacity=0.1}
%%		] {
%%			a_{11} & a_{12} & \cdots & a_{1n} \\
%%			a_{21} & a_{22} & \cdots & a_{2n} \\
%%			a & a & a & a \\
%%			a_{m1} & a_{m2} & \cdots & a_{mn} \\
%%		};
%%		\matrix (M) [matrix of math nodes, below=0.5cm of title1,
%%		left delimiter=(, right delimiter=),
%%		column sep=0.5em, row sep=0.5em,
%%		fill=cyan!2, draw=cyan!40, rounded corners=3pt, inner sep=10pt, drop shadow={opacity=0.1}] {
%%			|[entryNode, fill=cyan!5]| a_{11} & |[entryNode, fill=cyan!5]| a_{12} & |[dotsNode]| \cdots & |[entryNode, fill=cyan!5]| a_{1n} \\
%%			|[entryNode, fill=cyan!5]| a_{21} & |[entryNode, fill=cyan!5]| a_{22} & |[dotsNode]| \cdots & |[entryNode, fill=cyan!5]| a_{2n} \\
%%			|[dotsNode]| \vdots               & |[dotsNode]| \vdots               & |[dotsNode]| \ddots & |[dotsNode]| \vdots               \\
%%			|[entryNode, fill=cyan!5]| a_{m1} & |[entryNode, fill=cyan!5]| a_{m2} & |[dotsNode]| \cdots & |[entryNode, fill=cyan!5]| a_{mn} \\
%%		};
%		
%%		% Row/Col Highlight and Dimensions
%%		% Reverse coordinate order to flip braces instead of using 'mirror'
%%		\draw [revbrace, cyan!80] (M-1-4.north east) -- (M-1-1.north west) node [midway, above=12pt, font=\small\bfseries] {$n$ Entries per Row};
%%		\draw [revbrace, cyan!80] (M-4-1.south west) -- (M-1-1.north west) node [midway, left=12pt, font=\small\bfseries] {$m$ Rows};
%%		
%%		% Highlight specific entry link (e.g., a_12)
%%		\node[entryNode, fill=purple!30, draw=purple!80, thick] (highlightEntryM) at (M-1-2) {$a_{12}$};
%%		
%%		% ==============================================================================
%%		% PHASE 2: FLATTENING / INDEPENDENT ENTRIES (Purple Transition)
%%		% ==============================================================================
%%		\node[labelStyle, below=1.2cm of M, color=purple] (title2) {STEP 2: Map to $mn$ Independent Choices};
%%		
%%		% Flow arrows from matrix to flattening list
%%		\draw[arrowFlow, purple!80, dashed, curved arrow/.style={to path={(\tikztostart) .. controls ++(1,-1) and ++(-1,1) .. (\tikztotarget)}}]
%%		(M.south west) to[out=-90, in=90] ([xshift=-1cm]title2.north);
%%		\draw[arrowFlow, purple!80, dashed]
%%		(M.south east) to[out=-90, in=90] ([xshift=1cm]title2.north);
%%		
%%		% The Flattened List
%%		\matrix (Flattened) [matrix of math nodes, below=0.3cm of title2,
%%		nodes={entryNode, fill=purple!10, rounded corners=2pt},
%%		column sep=0.4em, inner sep=8pt, rounded corners=5pt,
%%		colorFlattenBg, drop shadow={opacity=0.1}] {
%%			a_{11} & a_{12} & \cdots & a_{1n} & a_{21} & \cdots & a_{mn} \\
%%		};
%%		
%%		% Track the highlighted entry
%%		\node[entryNode, fill=purple!30, draw=purple!80, thick, rounded corners=2pt] (highlightEntryF) at (Flattened-1-2) {$a_{12}$};
%%		\draw[->, purple, thick, shorten >=1pt, shorten <=1pt] (highlightEntryM) -- (highlightEntryF);
%%		
%%		% ==============================================================================
%%		% PHASE 3: APPLYING FIELD SIZE (Orange)
%%		% ==============================================================================
%%		\node[labelStyle, below=1.2cm of Flattened, color=orange!80!black] (title3) {STEP 3: Each Entry has $|\mathbb{F}_q|$ Possibilities};
%%		
%%		\draw[arrowFlow, orange!80] (title2.south) -- (title3.north);
%%		
%%		% The Choice Nodes (q's)
%%		\node[choiceButton, below=0.3cm of title3] (q1) {$q$};
%%		\node[right=1.2cm of q1] (dots1) {$\cdots$};
%%		\node[right=1.2cm of dots1] (qm) {$q$};
%%		\node[left=1.2cm of q1] (dots2) {$\cdots$};
%%		\node[left=1.2cm of dots2] (qmn) {$q$};
%%		
%%		% Map from entries to choices
%%		\draw[->, dashed, orange, thick] (Flattened-1-1.south) -- (qmn.north) node[midway, left, font=\scriptsize, text=orange!80!black] {choose};
%%		\draw[->, dashed, orange, thick] (highlightEntryF.south) -- (q1.north) node[midway, right, font=\scriptsize, text=purple] {choose};
%%		\draw[->, dashed, orange, thick] (Flattened-1-7.south) -- (qm.north) node[midway, right, font=\scriptsize, text=orange!80!black] {choose};
%%		
%%		% Label describing 'q' choice
%%		\node[right=1.5cm of Flattened-1-7, align=left, font=\scriptsize, text=orange!80!black] (chooseDesc) {Choices are independent:\\any $a_{ij} \in \{f_1, f_2, \dots, f_q\}$};
%%		\draw[->, thin, orange!50] (chooseDesc.west) -- ([xshift=10pt]Flattened-1-7.east);
%%		
%%		% ==============================================================================
%%		% PHASE 4: THE RESULT (Gold Plaque)
%%		% ==============================================================================
%%		% Reverse brace order to put on bottom
%%		\draw [revbrace, orange!70, thick, decoration={brace, amplitude=12pt}] ([xshift=10pt]qm.south east) -- ([xshift=-10pt]qmn.south west)
%%		node [midway, below=25pt, labelStyle, color=teal, font=\bfseries\large] (prodRule) {Product Rule Application};
%%		
%%		\node[resultPlaque, below=0.8cm of prodRule] (conclusion) {
%%			\LARGE \textbf{Total Possible Matrices:} \\
%%			\vspace{0.2cm}
%%			\begin{minipage}{0.8\linewidth}
%%				\centering
%%				\color{orange!60!black}
%%				$\underbrace{q \times q \times \dots \times q}_{mn \text{ factors}}$ = \Huge \textcolor{black}{$q^{mn}$} \\[0.3cm]
%%				\large \textbf{Result:} \fbox{\large $\bigl|\mathbb{F}_q^{m\times n}\bigr| = q^{mn}$} \hspace{0.5cm} \fbox{\large $\bigl|\mathbb{F}_q^{m\times m}\bigr| = q^{m^2}$}
%%			\end{minipage}
%%		};
%	\end{tikzpicture}
%\end{document}

\documentclass[tikz,border=20pt]{standalone}
\usepackage{amsmath, amssymb}
\usetikzlibrary{matrix, positioning, decorations.pathreplacing, shadings, shadows, arrows.meta, calc}

\begin{document}
	\begin{tikzpicture}[
		>=Stealth,
		node distance=1.5cm,
		% Define distinctive colors
		colorMatrixBg/.style={cyan!10, draw=cyan!60, thick},
		colorFlattenBg/.style={purple!15, draw=purple!60, thick},
		colorChoicesBg/.style={orange!20, draw=orange!70, thick},
		colorFinalBg/.style={yellow!25, draw=yellow!80, thick, drop shadow},
		% Modern styles for nodes
		entryNode/.style={draw, minimum size=2.2em, inner sep=0pt, font=\small, anchor=center, rounded corners=2pt},
		dotsNode/.style={minimum size=2.2em, inner sep=0pt, anchor=center, draw=none, fill=none},
		choiceButton/.style={circle, minimum size=2.5em, colorChoicesBg, font=\bfseries\large, drop shadow={opacity=0.2}},
		resultPlaque/.style={rectangle, rounded corners=8pt, colorFinalBg, inner sep=15pt, align=center},
		labelStyle/.style={font=\sffamily\bfseries\color{darkgray}, align=center},
		arrowFlow/.style={->, thick, shorten >=2pt, shorten <=2pt},
		% Corrected brace style
		revbrace/.style={decorate, decoration={brace, amplitude=8pt, raise=6pt}}
		]
		
		% ==============================================================================
		% PHASE 1: THE INPUT MATRIX (Cyan)
		% ==============================================================================
		\node[labelStyle] (title1) {$A \in \mathbb{F}_q^{m \times n}$};
		
		% FIX: Replaced \matrix with \node[matrix] and added ampersand replacement
		\node (M) [matrix of math nodes, ampersand replacement=\&, right=1.5cm of title1,
		left delimiter=(, right delimiter=),
		column sep=0.5em, row sep=0.5em,
		%fill=cyan!2, 
		%draw=cyan!40, 
		%rounded corners=3pt, inner sep=10pt, drop shadow={opacity=0.1}
		] {
			|[entryNode, fill=cyan!5, draw=cyan!40]| a_{11} \& |[entryNode, fill=cyan!5, draw=cyan!40]| a_{12} \& |[dotsNode]| \cdots \& |[entryNode, fill=cyan!5, draw=cyan!40]| a_{1n} \\
			|[entryNode, fill=cyan!5, draw=cyan!40]| a_{21} \& |[entryNode, fill=cyan!5, draw=cyan!40]| a_{22} \& |[dotsNode]| \cdots \& |[entryNode, fill=cyan!5, draw=cyan!40]| a_{2n} \\
			|[dotsNode]| \vdots               \& |[dotsNode]| \vdots               \& |[dotsNode]| \ddots \& |[dotsNode]| \vdots               \\
			|[entryNode, fill=cyan!5, draw=cyan!40]| a_{m1} \& |[entryNode, fill=cyan!5, draw=cyan!40]| a_{m2} \& |[dotsNode]| \cdots \& |[entryNode, fill=cyan!5, draw=cyan!40]| a_{mn} \\
		};
		
		% Row/Col Highlight and Dimensions
		\draw [revbrace, cyan!80] (M-1-1.north west) -- (M-1-4.north east) node [midway, above=12pt, font=\small\bfseries] {$n$};
		\draw [revbrace, cyan!80] ([xshift=-10pt]M-4-1.south west) -- ([xshift=-10pt]M-1-1.north west) node [midway, left=12pt, font=\small\bfseries] {$m$};
		
		% Highlight specific entry link
		\node[entryNode, fill=purple!30, draw=purple!80, thick] (highlightEntryM) at (M-1-1) {$a_{11}$};
		
		% ==============================================================================
		% PHASE 2: FLATTENING / INDEPENDENT ENTRIES (Purple Transition)
		% ==============================================================================
		\node[labelStyle, below=.5cm of M, color=purple] (title2) {$mn$ Independent Choices};
		
%		% Flow arrows from matrix to flattening list
%		\draw[arrowFlow, purple!80, dashed]
%		(M.south west) to[out=-90, in=90] ([xshift=-1cm]title2.north);
%		\draw[arrowFlow, purple!80, dashed]
%		(M.south east) to[out=-90, in=90] ([xshift=1cm]title2.north);
		
		% FIX: Replaced \matrix with \node[matrix] and used \& 
		\node (Flattened) [matrix of math nodes, ampersand replacement=\&, below=0.3cm of title2,
		column sep=0.4em, inner sep=8pt, rounded corners=5pt,
		fill=purple!15, draw=purple!60, thick
%		colorFlattenBg
		%, drop shadow={opacity=0.1}
		] {
			|[entryNode, fill=purple!10]| a_{11} \& |[entryNode, fill=purple!10]| a_{12} \& |[dotsNode]| \cdots \& |[entryNode, fill=purple!10]| a_{1n} \& |[entryNode, fill=purple!10]| a_{21} \& |[dotsNode]| \cdots \& |[entryNode, fill=purple!10]| a_{mn} \\
		};
		
		% Track the highlighted entry
		\node[entryNode, fill=purple!30, draw=purple!80, thick] (highlightEntryF) at (Flattened-1-1) {$a_{12}$};
		\draw[->, purple, thick, shorten >=1pt, shorten <=1pt, bend angle=45] (highlightEntryM.west) -| (highlightEntryF.north);
		
		% ==============================================================================
		% PHASE 3: APPLYING FIELD SIZE (Orange)
		% ==============================================================================
		\node[labelStyle, below=1.2cm of Flattened, color=orange!80!black] (title3) {STEP 3: Each Entry has $|\mathbb{F}_q|$ Possibilities};
		
		\draw[arrowFlow, orange!80] (title2.south) -- (title3.north);
		
		% The Choice Nodes (q's)
		\node[choiceButton, below=0.3cm of title3] (q1) {$q$};
		\node[right=1.2cm of q1] (dots1) {$\cdots$};
		\node[right=1.2cm of dots1] (qm) {$q$};
		\node[left=1.2cm of q1] (dots2) {$\cdots$};
		\node[left=1.2cm of dots2] (qmn) {$q$};
		
		% Map from entries to choices
		\draw[->, dashed, orange, thick] (Flattened-1-1.south) -- (qmn.north) node[midway, left, font=\scriptsize, text=orange!80!black] {choose};
		\draw[->, dashed, orange, thick] (highlightEntryF.south) -- (q1.north) node[midway, right, font=\scriptsize, text=purple] {choose};
		\draw[->, dashed, orange, thick] (Flattened-1-7.south) -- (qm.north) node[midway, right, font=\scriptsize, text=orange!80!black] {choose};
		
%		% Label describing 'q' choice
%		\node[right=1.5cm of Flattened-1-7, align=left, font=\scriptsize, text=orange!80!black] (chooseDesc) {Choices are independent:\\any $a_{ij} \in \mathbb{F}_q$};
%		\draw[->, thin, orange!50] (chooseDesc.west) -- ([xshift=10pt]Flattened-1-7.east);
%		
%		% ==============================================================================
%		% PHASE 4: THE RESULT (Gold Plaque)
%		% ==============================================================================
%		\draw [revbrace, orange!70, thick, decoration={brace, amplitude=12pt}] ([xshift=10pt]qm.south east) -- ([xshift=-10pt]qmn.south west)
%		node [midway, below=25pt, labelStyle, color=teal, font=\bfseries\large] (prodRule) {Product Rule Application};
%		
%		\node[resultPlaque, below=0.8cm of prodRule] (conclusion) {
%			\LARGE \textbf{Total Possible Matrices:} \\[0.2cm]
%			\color{orange!60!black}
%			$\underbrace{q \times q \times \dots \times q}_{mn \text{ factors}} =$ {\Huge \textcolor{black}{$q^{mn}$}} \\[0.4cm]
%			\large \textbf{Result:} \fbox{\large $\bigl|\mathbb{F}_q^{m\times n}\bigr| = q^{mn}$} \hspace{0.5cm} \fbox{\large $\bigl|\mathbb{F}_q^{m\times m}\bigr| = q^{m^2}$}
%		};
%		
	\end{tikzpicture}
\end{document}